\documentclass[12pt,a4paper]{ctexart}

% ── 编码与字体 ──
\usepackage{amsmath,amssymb}

% ── 页面设置 ──
\usepackage[margin=2.5cm]{geometry}
\usepackage{setspace}
\onehalfspacing

% ── 插图与表格 ──
\usepackage{graphicx}
\usepackage{float}
\usepackage{caption}
\usepackage{booktabs}
\usepackage{array}

% ── 颜色与链接 ──
\usepackage[colorlinks=true,linkcolor=blue,urlcolor=blue]{hyperref}

% ── 代码高亮 ──
\usepackage{listings}
\usepackage{xcolor}

\lstset{
    language=C++,
    basicstyle=\small\ttfamily,
    backgroundcolor=\color{gray!5},
    frame=single,
    breaklines=true,
    numbers=left,
    numberstyle=\tiny,
    keywordstyle=\color{blue},
    commentstyle=\color{green!50!black},
    stringstyle=\color{red!50!orange},
    tabsize=4,
    captionpos=b
}

\title{\vspace{-2cm}桌面宠物 --- 设计文档}
\author{}
\date{\today}

\begin{document}

\maketitle
\thispagestyle{empty}

\begin{abstract}
本文档介绍了一个基于 Windows GDI+ 与 WebView2 的桌面宠物程序的设计与实现。
该程序使用分层窗口（Layered Window）技术渲染带 Alpha 通道的 PNG 图片，
在屏幕上显示一只可交互的宠物，支持拖拽、点击反馈、物理下落、地面漫游以及 AI 聊天等功能。
\end{abstract}

\tableofcontents
\newpage

% ═══════════════════════════════════════════
\section{项目概述}
% ═══════════════════════════════════════════

\subsection{功能特性}
\begin{itemize}
    \item \textbf{拟物物理}：宠物受简单重力影响可下落，碰到地面后在地面上左右漫游。
    \item \textbf{拖拽交互}：用户可拖拽宠物，松开后根据拖拽力度和方向产生初速度。
    \item \textbf{点击反馈}：单击宠物时按权重随机切换显示不同图片。
    \item \textbf{方向感知}：行走、拖拽、下落等状态均区分左/右方向，自动选择对应文件夹 (\texttt{photo\_left/}、\texttt{photo\_right/}) 中的图片。
    \item \textbf{AI 聊天}：双击宠物打开聊天窗口，基于 WebView2 加载 HTML 聊天界面\cite{MidsummerBird}，通过 JavaScript 与 AI 后端通信。
    \item \textbf{可视化配置}：右键弹出设置窗口，可实时调整各图片的随机权重以及聊天颜色。
    \item \textbf{系统托盘}：支持最小化到托盘，双击托盘图标切换显示/隐藏。
\end{itemize}

\subsection{开发环境}
\begin{itemize}
    \item 编译器：MinGW-w64 (GCC) 8.1.0，使用 g++ 编译
    \item 图形库：GDI+（Windows GDI Plus）
    \item Web 引擎：Microsoft WebView2 (Edge Chromium)
    \item 构建方式：命令行编译 / VS Code 任务
    \item 源码行数：约 2,400 行
\end{itemize}

% ═══════════════════════════════════════════
\section{系统架构}
% ═══════════════════════════════════════════

程序整体采用单一 C 源文件（实际使用 C++ 编译），以 WinMain 为入口，围绕 Windows 消息循环构建事件驱动架构。模块划分如下：

\begin{figure}[H]
\centering
\begin{minipage}{0.9\textwidth}
\begin{lstlisting}[frame=single,numbers=none]
WinMain
  |-- GDI+ 初始化
  |-- 图片资源加载
  |-- 窗口创建 (WS_EX_LAYERED | WS_EX_TOPMOST)
  |-- 消息循环
       |-- WndProc (主窗口)
       |    |-- WM_TIMER  -> StepPet + RenderPet
       |    |-- WM_LBUTTONDOWN/UP/MOVE -> 点击/拖拽
       |    |-- WM_LBUTTONDBLCLK -> 聊天窗口
       |    |-- WM_RBUTTONUP -> 设置窗口
       |    |-- WM_WINDOWPOSCHANGING -> 强制置顶
       |    +-- WM_TRAYICON -> 托盘操作
       |
       |-- ChatWndProc (聊天窗口)
       |    |-- WM_CREATE -> 创建子控件
       |    |-- WM_COMMAND -> 发送/配置按钮
       |    +-- WM_CTLCOLOR -> 自定义颜色
       |
       |-- WebViewHostWndProc (WebView 容器)
       |    |-- WM_SIZE -> 调整 WebView 尺寸
       |    +-- WM_CLOSE -> 隐藏
       |
       +-- SettingsWndProc (设置窗口)
            +-- WM_COMMAND -> 应用/关闭/颜色选择
\end{lstlisting}
\end{minipage}
\end{figure}

% ═══════════════════════════════════════════
\section{核心技术}
% ═══════════════════════════════════════════

\subsection{分层窗口渲染}

宠物窗口使用 \texttt{WS\_EX\_LAYERED} 扩展样式创建，通过 \texttt{UpdateLayeredWindow} API
将包含 Alpha 通道的 32 位 DIB (Device Independent Bitmap) 直接绘制到屏幕上。
配合 \texttt{WS\_EX\_TOPMOST} 实现置顶显示。

渲染流程如代码 \ref{lst:render} 所示：创建内存 DC $\to$ 创建 DIBSection $\to$
使用 GDI+ Graphics 绘制 PNG $\to$ 调用 \texttt{UpdateLayeredWindow}。

\begin{lstlisting}[caption={RenderPet 核心渲染流程},label={lst:render},language=C++]
static void RenderPet(HWND hwnd) {
    // 1. 创建屏幕兼容 DC 和 32-bit DIBSection
    HDC screenDC = GetDC(NULL);
    HDC memDC = CreateCompatibleDC(screenDC);
    BITMAPINFO bmi = {};
    bmi.bmiHeader.biWidth  = g_pet.width;
    bmi.bmiHeader.biHeight = -g_pet.height; // 自上而下
    bmi.bmiHeader.biPlanes = 1;
    bmi.bmiHeader.biBitCount = 32;
    bmi.bmiHeader.biCompression = BI_RGB;
    HBITMAP dib = CreateDIBSection(screenDC, &bmi, DIB_RGB_COLORS, &bits,...);

    // 2. 用 GDI+ 绘制 PNG（支持透明通道）
    Graphics graphics(memDC);
    graphics.SetCompositingMode(CompositingModeSourceCopy);
    graphics.Clear(Color(0,0,0,0));          // 全透明底色
    graphics.SetCompositingMode(CompositingModeSourceOver);
    graphics.DrawImage(g_image, 0, 0, w, h); // 叠放 PNG

    // 3. 更新分层窗口
    BLENDFUNCTION bf = { AC_SRC_OVER, 0, 255, AC_SRC_ALPHA };
    UpdateLayeredWindow(hwnd, screenDC, &pos, &size, memDC, &src, 0, &bf, ULW_ALPHA);
}
\end{lstlisting}

\subsection{物理与状态机}

宠物状态分为 \textbf{飞行} (airborne) 和 \textbf{地面} (ground) 两种。
每帧 (\texttt{WM\_TIMER}, 10ms) 执行 \texttt{StepPet}，流程如下：

\begin{enumerate}
    \item 更新位置：$x \gets x + v_x$，$y \gets y + v_y$
    \item 边界检查：碰到左右墙反弹（$v_x \gets -v_x$），触地则 $v_y = 0$，进入地面模式
    \item 地面任务：接地后随机选择一个地面任务（定时左移/右移/到边缘/停止）
    \item 图片切换：根据方向（$v_x$ 符号）和状态选择对应的 PNG
\end{enumerate}

\subsection{WebView2 集成}

聊天功能通过 WebView2 加载本地 HTML 文件实现\cite{MidsummerBird}。通信采用双向 JSON 消息传递：

\begin{itemize}
    \item C++ $\to$ JS：\texttt{PostWebMessageAsJson} 发送包含用户文本的 JSON
    \item JS $\to$ C++：通过 \texttt{webMessageReceived} 事件接收 AI 回复
\end{itemize}

采用 COM 回调模式，定义了三个回调类：
\begin{itemize}
    \item \texttt{WebViewEnvCompletedHandler} --- 环境创建完成
    \item \texttt{WebViewControllerCompletedHandler} --- 控制器创建完成
    \item \texttt{WebViewMessageReceivedHandler} --- 接收 JS 消息
\end{itemize}

\subsection{窗口置顶优先级}

为解决宠物窗口被任务栏遮挡的问题，采用双层保障策略：

\begin{lstlisting}[caption={强制置顶实现},language=C++]
// 策略 1: 拦截 Z-order 变化
case WM_WINDOWPOSCHANGING: {
    WINDOWPOS* wp = (WINDOWPOS*)lParam;
    wp->hwndInsertAfter = HWND_TOPMOST;
    wp->flags &= ~SWP_NOZORDER;
    return 0;
}

// 策略 2: 定时器轮询（每 500ms）
case WM_TIMER:
    // ... 更新物理与渲染 ...
    if ((GetTickCount() / 500) != g_lastTopmostCheck) {
        g_lastTopmostCheck = GetTickCount() / 500;
        SetWindowPos(hwnd, HWND_TOPMOST, 0, 0, 0, 0,
                     SWP_NOMOVE | SWP_NOSIZE | SWP_NOACTIVATE);
    }
\end{lstlisting}

% ═══════════════════════════════════════════
\section{图片资源管理}
% ═══════════════════════════════════════════

\subsection{目录结构}

程序从 \texttt{photo\_left/} 和 \texttt{photo\_right/} 两个目录读取图片资源，
基于当前移动方向自动选择：

\begin{verbatim}
Desktop pet/
  |-- text1.c          # 源码
  |-- output/text1.exe # 编译输出
  |-- photo_left/      # 向左移动时使用的图片
  |     |-- 0.png      #  行走
  |     |-- 1.png      #  停止触发（权重可配）
  |     |-- 2.png      #  低高度下落 / 落地
  |     |-- 3.png      #  拖拽
  |     |-- 4.png      #  停止触发（权重可配）
  |     |-- 5.png      #  下落
  |     |-- 6.png      #  点击触发（权重可配）
  |     |-- 7.png      #  点击触发（权重可配）
  |     +-- 8.png      #  落地特殊（权重可配）
  |-- photo_right/     # 向右移动时使用的图片
  |     +-- (同上)
  |-- ai_chat/         # AI 聊天界面
  |     +-- Midsummer's Bird.html
  +-- external/        # WebView2 SDK
\end{verbatim}

\subsection{权重随机机制}

程序通过加权随机算法实现多样化反馈。以点击事件为例：

\begin{equation}
\text{选中图片 i 的概率} = \frac{w_i}{\sum_{j} w_j}
\end{equation}

其中 $w_i$ 为第 $i$ 张图片的权重，所有有效图片的权重和作为分母。
各权重可在设置窗口中实时调整。代码 \ref{lst:weight} 展示了实现：

\begin{lstlisting}[caption={加权随机选择},label={lst:weight},language=C++]
static Image* PickClickImageByDirection(BOOL toLeft) {
    int total = 0;
    for (int i = 0; i < kClickItemCount; ++i) {
        Image* img = toLeft ? g_clickItems[i].left : g_clickItems[i].right;
        if (img && g_clickItems[i].weight > 0)
            total += g_clickItems[i].weight;
    }
    if (total <= 0) return NULL;
    int r = rand() % total;
    for (int i = 0; i < kClickItemCount; ++i) {
        Image* img = toLeft ? g_clickItems[i].left : g_clickItems[i].right;
        if (img && g_clickItems[i].weight > 0) {
            if (r < g_clickItems[i].weight) return img;
            r -= g_clickItems[i].weight;
        }
    }
    return NULL;
}
\end{lstlisting}

% ═══════════════════════════════════════════
\section{编译与运行}
% ═══════════════════════════════════════════

\subsection{前置条件}
\begin{itemize}
    \item MinGW-w64 (g++) 或任何 C++17 编译器
    \item WebView2 Runtime（Windows 11 自带，Windows 10 需安装）
    \item GDI+ SDK（MinGW 自带）
\end{itemize}

\subsection{编译命令}
\begin{lstlisting}[language=bash,numbers=none]
g++ -I "external/webview2/pkg/build/native/include" ^
    -g3 -mwindows text1.c -o output/text1.exe ^
    -lgdi32 -luser32 -lmsimg32 -lgdiplus -lole32 -luuid -lshell32
\end{lstlisting}

\subsection{运行}
双击 \texttt{output/text1.exe} 或从命令行启动。

% ═══════════════════════════════════════════
\section{核心数据结构}
% ═══════════════════════════════════════════

\subsection{PetState --- 宠物物理状态}
\begin{lstlisting}[language=C++]
typedef struct {
    int x, y;       // 屏幕坐标（像素）
    int vx, vy;     // 速度（像素/帧，每帧约 16ms）
    int width, height; // 窗口尺寸
} PetState;
\end{lstlisting}

\subsection{GroundTask --- 地面行为枚举}
\begin{lstlisting}[language=C++]
typedef enum {
    TASK_NONE,
    TASK_MOVE_LEFT_EDGE,   // 向左走到边缘
    TASK_MOVE_RIGHT_EDGE,  // 向右走到边缘
    TASK_MOVE_LEFT_TIME,   // 向左走固定时长 (5s)
    TASK_MOVE_RIGHT_TIME,  // 向右走固定时长 (5s)
    TASK_STOP_TIME         // 停止 (1.5s)，仅在上一个任务为
                           // MOVE_LEFT_TIME/RIGHT_TIME 后可触发
} GroundTask;
\end{lstlisting}

\subsection{ClickImageItem --- 可配置的图片资源}
\begin{lstlisting}[language=C++]
typedef struct {
    const wchar_t* fileName; // 图片文件名
    int weight;              // 随机权重
    Image* left;             // 左方向时使用的 GDI+ Image
    Image* right;            // 右方向时使用的 GDI+ Image
} ClickImageItem;
\end{lstlisting}

% ═══════════════════════════════════════════
\section{改进与优化}
% ═══════════════════════════════════════════

本代码经过以下优化处理：

\begin{enumerate}
    \item \textbf{清理重复清理代码}：WinMain 末尾的图片资源释放原先重复 38 行，
    通过 \texttt{SAFE\_DELETE\_IMAGE} 宏压缩为 14 行。
    \item \textbf{消除编译警告}：将结构体初始化 \texttt{\{ 0 \}} 统一改为 \texttt{\{\}}，
    消除所有 "-Wmissing-field-initializers" 警告；移除未使用的变量 \texttt{g\_trayHwnd}。
    \item \textbf{强制置顶}：添加 WM\_WINDOWPOSCHANGING 拦截 + 定时器轮询，
    双重保障宠物不被任务栏遮挡。
    \item \textbf{使用 C++ 编译器}：原代码使用 GDI+ 和 C++ 特性（class、new/delete、reinterpret\_cast），
    必须用 g++ 而非 gcc 编译。
\end{enumerate}

% ═══════════════════════════════════════════
\begin{thebibliography}{99}
% ═══════════════════════════════════════════

\bibitem{MidsummerBird}
Soaring Bird. \textit{Midsummer's Bird --- An Elegant, Powerful, and Fully Localized AI Role-Playing Chat Frontend}. \url{https://github.com/SoaringBird/Midsummer-s-Bird}, 2025. MIT License.

\bibitem{SillyTavern}
SillyTavern. \textit{SillyTavern --- LLM Frontend for Power Users}. \url{https://github.com/SillyTavern/SillyTavern}.

\bibitem{WebView2}
Microsoft. \textit{WebView2 --- Microsoft Edge WebView2 Control}. \url{https://learn.microsoft.com/en-us/microsoft-edge/webview2/}, 2025.

\bibitem{MinGW}
MinGW-w64. \textit{MinGW-w64 --- GCC for Windows 64 \& 32-bit}. \url{https://www.mingw-w64.org/}.

\end{thebibliography}

\end{document}
